\originalchapter{幂级数与恒等定理}
\sourceinfo{18.04 Topic 7, pp.1--10及未完成的收敛附录; 18.112 L4. 一致收敛估计为编者补充}
\originalsection{收敛半径}
\begin{definition}
级数 $\sum c_n$ 收敛, 指部分和趋于有限复数. 绝对收敛指 $\sum|c_n|<\infty$. 函数级数在集合 $K$ 上一致收敛, 指尾和可由同一个与点无关的界控制. 在每个紧集上一致收敛称为局部一致收敛.
\end{definition}
绝对收敛保证 Cauchy 条件: 任何足够后的尾和模不超过绝对值尾和. 复数完备性来自两个实坐标的完备性. 若 $|f_n(z)|\le M_n$ 于 $K$ 且 $\sum M_n<\infty$, 则尾和一致趋零; 这是 Weierstrass 判别法. 一致收敛使连续性和有限长度曲线上的积分可以交换极限, 因为积分误差至多为长度乘上一致误差.

比值判别中若 $|c_{n+1}/c_n|\to q<1$, 则后续项由几何级数控制, 绝对收敛; 若 $q>1$, 项不趋零, 级数发散. 根值判别使用 $\limsup |c_n|^{1/n}$, 无需极限存在. 值等于1时两者都不作结论.
\begin{theorem}[收敛半径]\label{thm:radius}
设 $\sum_{n\ge0}a_n(z-a)^n$ 为幂级数, 令 $R^{-1}=\limsup_{n\to\infty}|a_n|^{1/n}$, 按 $1/0=\infty$, $1/\infty=0$ 解释. 级数在 $|z-a|<R$ 绝对且局部一致收敛, 在 $|z-a|>R$ 发散. $R>0$ 时其和全纯, 可逐项求导、积分, 导数级数仍有半径 $R$.
\end{theorem}
\begin{proof}
对 $r<R$, 选 $r<\rho<R$. 由上极限, 对充分大 $n$ 有 $|a_n|\le \rho^{-n}$. 在 $|z-a|\le r$ 上项由 $(r/\rho)^n$ 控制, 所以一致绝对收敛. 若 $|z-a|>R$, 选择 $R<\rho<|z-a|$, 则有无穷多个 $n$ 满足 $|a_n|>\rho^{-n}$, 对应项模不趋零. 端点 $R=0,\infty$ 按同一论证处理.

部分和都是全纯多项式, 局部一致极限定理给和全纯、导数收敛. 导数系数 $na_n$ 的 $n$ 次根多一个趋于1的因子, 移动一个指标不改变半径; 也可用 $n(r/\rho)^n$ 的可求和性直接得到导数一致收敛. 对任意固定分段光滑路径, 其像是收敛圆盘内的紧集, 积分误差估计允许逐项积分.
\end{proof}
边界需要另查. $\sum z^n$ 在单位圆周处处发散; $\sum z^n/n^2$ 在闭单位圆盘上一致绝对收敛. 两者半径均为1. 不可把圆内局部一致收敛误说成在整个开圆盘上一致收敛.

\originalsection{Taylor 展开}
\begin{theorem}\label{thm:taylor}
若 $f$ 在 $D(a,R)$ 全纯, 则
\[
f(z)=\sum_{n=0}^\infty\frac{f^{(n)}(a)}{n!}(z-a)^n,\qquad |z-a|<R.
\]
展开在每个较小闭圆盘上一致绝对收敛, 且系数唯一.
\end{theorem}
\begin{proof}
取 $r<\rho<R$, 对 $|z-a|\le r$, $|\zeta-a|=\rho$, 几何级数给
\[
\frac1{\zeta-z}=\sum_{n=0}^N\frac{(z-a)^n}{(\zeta-a)^{n+1}}
 +\frac{(z-a)^{N+1}}{(\zeta-a)^{N+1}(\zeta-z)}.
\]
代入 Cauchy 公式, 系数积分由导数公式等于 $f^{(n)}(a)/n!$. 余项模不超过 $M_\rho(r/\rho)^{N+1}\rho/(\rho-r)$, 一致趋零. 绝对收敛由 Cauchy 系数估计得到. 任意幂级数表达可以逐项求导并在 $a$ 取值, 必须具有上述系数, 所以唯一.
\end{proof}
该证明只依赖 Cauchy 公式与几何级数, 并未使用最大模原理. 它补足第4章从局部常数到整个域恒定的最后环节, 因而不存在循环依赖. 全纯与解析的等价也至此建立: Taylor 定理给一个方向, 幂级数和的全纯性给反方向.

几个常用展开为
\begin{align*}
\ee^z&=\sum_{n\ge0}\frac{z^n}{n!},&
\sin z&=\sum_{n\ge0}\frac{(-1)^nz^{2n+1}}{(2n+1)!},\\
\cos z&=\sum_{n\ge0}\frac{(-1)^nz^{2n}}{(2n)!},&
\Log(1+z)&=\sum_{n\ge1}\frac{(-1)^{n-1}z^n}{n},\quad |z|<1.
\end{align*}
最后一式可从 $1/(1+z)$ 的几何级数逐项积分, 用 $\Log1=0$ 确定常数. 所使用的分支在整个圆盘内全纯. 展开中心改变时必须重新写成 $z-a$ 的幂; 如 $1/(1-z)$ 在 $a=5$ 的展开为 $-\tfrac14\sum_{n\ge0}(-(z-5)/4)^n$, 半径4.

\begin{example}
求 $\ee^z/(1-z)$ 在0的 Taylor 系数和半径.
\end{example}
\begin{solution}
在 $|z|<1$ 内两个级数绝对收敛, 其乘积可按总次数收集. 系数为 $a_n=\sum_{k=0}^n1/k!$. 因此展开是 $\sum a_nz^n$. $a_n\to\ee\ne0$, 根值判别给半径1. 从函数看, 点1是不能消去的极点; 若半径大于1, 幂级数和会在1附近全纯且与原函数在交域相同, 矛盾.
\end{solution}

\originalsection{零点与恒等}
\begin{definition}
若 $f(a)=0$ 且存在最小正整数 $m$ 使 $f^{(m)}(a)\ne0$, 则零点阶数为 $m$. 此时 $f(z)=(z-a)^mg(z)$, $g$ 在 $a$ 附近全纯, $g(a)\ne0$.
\end{definition}
\begin{theorem}[恒等定理]\label{thm:identity}
域上的全纯函数若零点在域内有聚点, 则恒为零. 因而非常零全纯函数的零点都孤立. 两个全纯函数若在域内有聚点的集合上相等, 则在整个域相等.
\end{theorem}
\begin{proof}
在零点聚点 $a$ 的 Taylor 展开若有首个非零系数, 则可分解为 $(z-a)^mg(z)$, 且小邻域内 $g\ne0$, 所以 $a$ 是孤立零点, 矛盾. 因而全部系数为零, $f$ 在 $a$ 邻域恒为零.

令 $S$ 为 $f$ 在邻域恒为零的点集. $S$ 开且非空. 若 $b$ 是 $S$ 在域内的极限点, 所有阶导数在 $S$ 为零, 由连续性在 $b$ 也为零, Taylor 展开使 $b\in S$. 所以 $S$ 在域内又闭, 连通性迫使它是全域. 最后应用于两函数之差.
\end{proof}
聚点必须位于域内. $\sin(\pi/z)$ 在穿孔平面全纯, 零点 $1/n$ 聚到缺失的0, 并不恒为零. 连通性也不能删去: 在上下半平面分别定义不同常数, 可使两个全纯函数只在一个分支相等.

若 $f$ 在0的全部导数为零, 则全纯性迫使它局部恒为零. 这与实光滑函数不同: 实函数 $\ee^{-1/x^2}$ 在 $x\ne0$ 定义、在0补零后无限次可微, 零点所有导数为零却不恒为零. 决定复函数的刚性的是复可微和 Cauchy 理论, 不是仅仅导数的数量.

\originalsection{幂级数解微分方程}
解析系数微分方程可借未知 Taylor 系数的递推求解, 但形式递推后还须验证收敛. 以 $f''+f=0$, $f(0)=A,f'(0)=B$ 为例, 写 $f=\sum a_nz^n$, 逐次系数给 $(n+2)(n+1)a_{n+2}+a_n=0$. 初值使 $a_0=A,a_1=B$, 因而
\[
a_{2n}=\frac{(-1)^nA}{(2n)!},\qquad
 a_{2n+1}=\frac{(-1)^nB}{(2n+1)!}.
\]
所得两列半径无穷, 逐项微分合法, 所以 $f=A\cos z+B\sin z$ 确实是整解. 任意全纯解都有 Taylor 展开并满足同一递推, 所以在0附近唯一, 恒等定理使其在连通共同域唯一. 若方程在某点的最高导数系数为零, 该点递推可能失效, 不能把这个正常点例子的结论无条件推广到奇点.

例如 $f'=(1-z)^{-1}f$, $f(0)=1$, 在 $|z|<1$ 有系数关系 $(n+1)a_{n+1}=\sum_{k=0}^na_k$. 由归纳 $a_n=1$, 和为 $1/(1-z)$, 在0的收敛半径1恰受方程系数的奇点限制. 形式系数并不保证整平面收敛, 需像这里一样另作半径验证.

\originalsection{习题与解答}
\begin{exercise}
求 $\sum_{n\ge0}n!z^n$ 和 $\sum_{n\ge0}z^n/(n!)^2$ 的收敛半径.
\end{exercise}
\begin{solution}
第一列相邻项模比为 $(n+1)|z|$, 任意 $z\ne0$ 时最终大于1, 所以半径0. 第二列比值为 $|z|/(n+1)^2\to0$, 任意 $z$ 都绝对收敛, 半径无穷.
\end{solution}
\begin{exercise}
证明整函数若在所有实数上等于 $\cos x$, 则等于 $\cos z$ 于全平面; 若仅在所有整数上等于 $\cos n$, 能否得到同一结论?
\end{exercise}
\begin{solution}
实轴具有复平面内的聚点, 恒等定理适用. 整数集合没有有限聚点, $\cos z+\sin(\pi z)$ 是反例.
\end{solution}
\begin{exercise}
求 $(\sin z-z)/z^3$ 在0的连续补值与前两项 Taylor 展开.
\end{exercise}
\begin{solution}
正弦级数给函数为 $-1/6+z^2/120+O(z^4)$. 补值为 $-1/6$, 补后在整个平面全纯; 除法造成的零分母已经被三阶零点消去.
\end{solution}
